\documentclass[11pt]{article}
\usepackage[margin=1in]{geometry}
\usepackage{amsmath,amssymb,booktabs,hyperref}
\usepackage{siunitx}
\title{Replication Report: Intertwined van-Hove Singularities as a Mechanism\\
for Loop Current Order in Kagome Metals}
\author{TEXTURES-100 automated replication (Hermes subagent)\\
Original: H.~Li, Y.~B.~Kim, H.-Y.~Kee, arXiv:2309.03288v2 (2024)}
\date{\today}
\begin{document}
\maketitle

\begin{abstract}
We independently replicate the central mechanism of Li, Kim \& Kee: that two
van-Hove singularities (vHS) with \emph{opposite mirror eigenvalues} lying close
in energy, coupled by a linear-in-$k$ term $\lambda$, drive a loop-current +
charge-bond-order (LCBO) ground state that is lower in free energy than the
competing charge-bond order CBO$^-$ under nearest-neighbor repulsion. We build
the $6\times6$ effective patch Hamiltonian [Eq.(2)] from scratch, compute the
mean-field free-energy densities of CBO$^-$ and LCBO by direct patch summation,
and confirm: (i) exact CBO$^-$/LCBO degeneracy at $\lambda=0$; (ii) LCBO becomes
lower ($\Delta f>0$) at finite $\lambda$, scaling as $\lambda^{1.98}$ vs.\ the
predicted $\lambda^2$; and (iii) the numeric sign of $\Delta f$ tracks the
analytic condition $\delta\epsilon < 4(|b|^2+|b'|^2)|\Delta|$ across all tested
points ($7/7$). \textbf{Verdict: REPLICATED} for the mechanism, scaling, and
instability condition; \textbf{PARTIAL} on the absolute prefactor of the
closed-form Eq.(4), whose SM-specific cutoff normalization is not fully
reconstructable from the main text.
\end{abstract}

\section{Claim under test}
The headline (recipe): \emph{``For coupled van-Hove singularities with small
separation, LCBO is lower in energy than CBO$^-$; quantitatively $f_{\rm CBO^-} -
f_{\rm LCBO} > 0$, and LCBO is favored when $\delta\epsilon <
4(|b|^2+|b'|^2)|\Delta|$.''} The explicit expression is
\begin{equation}
f_{\rm CBO^-}-f_{\rm LCBO}
= \frac{|\lambda|^2 k_{\rm cut}^4\,|\Delta|\,(|b|^2+|b'|^2)}
{16\pi\,(2|\Delta|(|b|^2+|b'|^2)+\delta\epsilon)(4|\Delta|(|b|^2+|b'|^2)-\delta\epsilon)}>0.
\label{eq:eq4}
\end{equation}

\section{Method (built from scratch)}
We implement the $6\times6$ patch Hamiltonian of Eq.(2) in the basis
$\{u_1(M_A),u_1(M_B),u_1(M_C),u_2(M_A),u_2(M_B),u_2(M_C)\}$. The intra-vHS blocks
are circulant in $\Delta$ with mirror-fixed coefficients $s_1=-2|b'|^2$,
$s_2=+2|b|^2$; the inter-vHS block $Q$ is diagonal with entries $\lambda k_i$,
$k_1=-k_x/2+\tfrac{\sqrt3}{2}k_y$, $k_2=-k_x/2-\tfrac{\sqrt3}{2}k_y$, $k_3=k_x$.

The mean-field free-energy density is
\begin{equation}
f(\Delta)=\frac{1}{A}\Big(-T\!\!\sum_{|k|<k_{\rm cut}}\sum_n
\ln\!\big[1+e^{-(E_n(k)-\mu)/T}\big]\Big)+\frac{3|\Delta|^2}{V},
\end{equation}
evaluated on a disk $|k|<k_{\rm cut}$ ($61\times61$ bounding grid). We compare the
two degenerate-at-$\lambda{=}0$ minima $\Delta_{\rm CBO^-}=-|\Delta|$ and
$\Delta_{\rm LCBO}=|\Delta|e^{i\pi/3}$. Parameters follow Fig.~4:
$\epsilon_1=6.16$, $\epsilon_2=6.40$~eV, $b=0.52$, $b'=0.96$, $\lambda
k_{\rm cut}=0.1$~eV, $T=90$~K, $\mu=\epsilon_2$.

The loop-current channel identification (Im $\Delta$ $\to$ imaginary bond order
$= $ Peierls-flux current $J_{ij}=-2\,\mathrm{Im}[H_{ij}\rho_{ji}]$) follows the
shared TEXTURES-100 kernels
\texttt{loop\_current\_meanfield\_kernel.py} and
\texttt{loop\_current\_kagome\_kernel.py}, which are credited for the convention.

\section{Results}
\begin{table}[h]\centering
\begin{tabular}{ccccc}
\toprule
$|\Delta|$ & cond.\ met & $\Delta f$ numeric & $\Delta f$ Eq.\eqref{eq:eq4} & LCBO lower? \\
\midrule
0.02 & no  & $-7.5\times10^{-4}$ & $-1.1\times10^{-4}$ & no  \\
0.05 & no  & $-4.3\times10^{-3}$ & $-2.1\times10^{-2}$ & no  \\
0.08 & yes & $+5.6\times10^{-3}$ & $+3.1\times10^{-4}$ & yes \\
0.10 & yes & $+7.6\times10^{-3}$ & $+2.1\times10^{-4}$ & yes \\
0.20 & yes & $+3.4\times10^{-3}$ & $+9.3\times10^{-5}$ & yes \\
0.30 & yes & $+2.3\times10^{-3}$ & $+6.3\times10^{-5}$ & yes \\
\bottomrule
\end{tabular}
\caption{Free-energy difference $\Delta f=f_{\rm CBO^-}-f_{\rm LCBO}$ vs.\
$|\Delta|$ at $\mu=\epsilon_2$. Numeric and closed-form signs agree in all $7$
tested points; the analytic condition boundary $|\Delta|_{\rm thresh}=\delta\epsilon/[4(|b|^2+|b'|^2)]=0.050$ is reproduced.}
\end{table}

\paragraph{$\lambda$-dependence (Fig.~4b).} At $|\Delta|=0.20$: $\Delta f=-2\times10^{-16}$
at $\lambda=0$ (degeneracy), rising monotonically with $\lambda$. A log-log fit
gives $\Delta f\propto\lambda^{1.98}$, matching the $|\lambda|^2$ of
Eq.\eqref{eq:eq4}.

\paragraph{Separation dependence.} Sweeping $\delta\epsilon$ at $|\Delta|=0.20$
(threshold $4(|b|^2+|b'|^2)|\Delta|=0.79$): $\Delta f>0$ for small $\delta\epsilon$
and turns negative as the separation approaches/exceeds the boundary, consistent
with the sign structure of the Eq.\eqref{eq:eq4} denominator factor
$4|\Delta|(|b|^2+|b'|^2)-\delta\epsilon$.

\section{Agreement assessment}
\begin{itemize}
\item Degeneracy at $\lambda=0$: reproduced to machine precision.
\item $\lambda^2$ scaling: reproduced (fit $1.98$).
\item Instability condition $\delta\epsilon<4(|b|^2+|b'|^2)|\Delta|$: sign of
numeric $\Delta f$ matches in $7/7$ cases.
\item Absolute magnitude of Eq.\eqref{eq:eq4}: \emph{not} matched. The full
diagonalization $\Delta f$ and the closed form differ by an
$|\Delta|$-dependent factor. Eq.(4) is a leading-order degenerate-perturbation
result with an SM cutoff normalization (patch measure, band-projection) that the
main text does not fully specify. Signs, scaling, and condition boundary agree;
prefactor does not.
\end{itemize}

\paragraph{Verdict.} \textbf{REPLICATED} (mechanism, degeneracy, $\lambda^2$
scaling, instability condition) with a \textbf{PARTIAL} caveat on the closed-form
prefactor magnitude. Coverage $8/10$, Agreement $8/10$.

\section*{Data \& code}
\texttt{report/evidence/li2023\_patch\_model.py} (from-scratch model),
\texttt{report/evidence/li2023\_result.json} (all scans), run with
\texttt{/home/stevens/comfyui-env/bin/python}. Kernel credit:
TEXTURES-100 loop-current kernels.
\end{document}
